\par\bigskip\Needspace{9\baselineskip}
\pdfbookmark[2]{Zweite Methode, [5]}{ceres-second-5}
\begin{center}[5.]\end{center}
\begingroup\small
\noindent\hspace*{\parindent}[Nachdem auf diese Weise alle Glieder des Ausdrucks 2) auf die Form
\[
A\sin(iD\pm\text{\AstroSymbols ⚳}\mp\AscendingNode+kM-k'M')
\]
gebracht sind, wird das entsprechende
\[
\begin{aligned}
&5)\qquad\text{Glied in }\frac{\theta}{\sin i}=\\
&\quad-A\cos V\int\sin V.\sin(iD\pm\text{\AstroSymbols ⚳}\mp\AscendingNode+kM-k'M')\,n\,\mathrm{d}t\\
&\quad+A\sin V\int\cos V.\sin(iD\pm\text{\AstroSymbols ⚳}\mp\AscendingNode+kM-k'M')\,n\,\mathrm{d}t.
\end{aligned}
\]

Der nächste Schritt besteht darin, \(\displaystyle\genfrac{}{}{0pt}{}{\sin}{\cos}V\) nach Vielfachen von \(M\) zu entwickeln. Auch dies kann durch Interpolation geschehen. Man erhielt
\[
\begin{aligned}
\cos V&=0{,}07848+0{,}99308\cos M-0{,}07784\cos2M\quad\text{usw.}\\
\sin V&=\phantom{0{,}07848+}0{,}99461\sin M-0{,}07792\sin2M\quad\text{usw.}
\end{aligned}
\]

Nach dem Ausmultiplizieren mit diesen Größen unter den Integralzeichen in 5) lässt sich die Integration unmittelbar ausführen. Durch erneute Multiplikation mit \(\cos V\) beziehungsweise \(\sin V\) erhält man das entsprechende Glied in \(\theta\). Für das oben als Beispiel verwendete Glied 3) ergab sich so:
\SourceLeaf{413}
\[
\begin{aligned}
&-27''{,}47\sin(3D+\text{\AstroSymbols ⚳}-\AscendingNode)\\
&-\phantom{2}4''{,}56\sin(3D+\text{\AstroSymbols ⚳}-\AscendingNode+M)\\
&+\phantom{2}2''{,}26\sin(3D+\text{\AstroSymbols ⚳}-\AscendingNode-M)\\
&+\phantom{2}2''{,}81\sin(3D+\text{\AstroSymbols ⚳}-\AscendingNode+M').
\end{aligned}
\]

Die gesamten numerischen Ergebnisse finden sich im folgenden Brief. Daraus geht auch hervor, wie Glieder, deren Argumente sich nur um konstante Größen unterscheiden, zu einem einzigen zusammengefasst wurden.]
\par\endgroup
\begin{center}\rule{.30\linewidth}{.4pt}\end{center}
